\section{The two-layer architecture}
\label{sec:architecture}

\begin{lede}
Every mechanism in this document sits in one of two layers: an upper layer that
decides what each robot \emph{wants}, and a lower layer that decides what is
\emph{allowed}. Only the lower layer can prevent a collision, and nothing in
the upper layer is permitted to remove a move the lower layer would have
considered.
\end{lede}

TOLL-PIBT extends PIBT \citep{okumura2019pibt} --- a single-timestep,
decentralised multi-agent collision-avoidance algorithm --- with a high-level
layer that decides \emph{what} each robot wants (task, waypoint, preferred
direction, priority) without ever touching collision safety.

\begin{center}
\small\ttfamily
\begin{tabular}{@{}l@{}}
\toprule
\textbf{\textsf{HIGH LEVEL}} \\
task stream $\rightarrow$ assignment $\rightarrow$ waypoints $\rightarrow$ routes $\rightarrow$ directional demand \\
$\rightarrow$ aisle direction (hysteresis, drain) $\rightarrow$ reservations \\
$\rightarrow$ priorities $\rightarrow$ deadlock detection \& local recovery \\
\midrule
\multicolumn{1}{@{}c@{}}{\textsf{\itshape ranks and prices candidates; never removes one}} \\
\midrule
\textbf{\textsf{LOW LEVEL (PIBT)}} \\
candidates $\rightarrow$ safety rejection $\rightarrow$ score-sorted $\rightarrow$ priority \\
inheritance $\rightarrow$ backtracking $\rightarrow$ one synchronised timestep \\
\bottomrule
\end{tabular}
\end{center}

The invariant this document keeps returning to: \textbf{everything above the
line only reorders candidate moves; only the PIBT recursion
(\Cref{sec:pibt}) and its validation companions (\Cref{sec:pibt-conflicts})
decide collision-freedom.} No high-level formula can, by construction, produce
two robots on the same vertex.

The middle row of that diagram is where this system differs from its own
original specification, and from the usual way a directional rule is
implemented. The specification listed ``violates aisle direction'' and
``violates a reservation'' among the low level's hard rejection rules; the same
document's design principle said the high level never replaces PIBT's
collision checks or its backtracking. Those two cannot both hold, and
\Cref{sec:contribution-soft} works out which one had to give.

\Cref{tab:modules} maps the sections that follow onto the code.

\begin{table}[H]
\centering\small
\begin{tabular}{@{}lll@{}}
\toprule
\textbf{Module} & \textbf{Section} & \textbf{What it does} \\
\midrule
\code{types.py} & --- & enums, \code{Vertex}, \code{Reservation} \\
\code{config.py} & \ref{sec:parameters} & \code{Params}, ablation presets, factorial designs \\
\code{graph.py} & \ref{sec:graph} & grid, BFS distance maps, articulation points, bridges \\
\code{warehouse.py} & \ref{sec:graph} & map parsing, vertex metadata, aisle segmentation \\
\code{task.py} & \ref{sec:robots-tasks} & tasks, queue, four arrival processes \\
\code{robot.py} & \ref{sec:robots-tasks} & robot state \\
\code{congestion.py} & \ref{sec:congestion} & occupancy index, local / aisle / downstream congestion \\
\code{scoring.py} & \ref{sec:scoring} & proximity modes, weights, turning cost, $\score$ \\
\code{priority.py} & \ref{sec:priority} & task-class, waiting and blocked priority \\
\code{aisle\_manager.py} & \ref{sec:aisles} & demand, hysteresis, drain-before-reverse, reservations \\
\code{assignment.py} & \ref{sec:assignment} & greedy $J(i,\tau)$ matching, waypoint transitions \\
\code{routing.py} & \ref{sec:routing} & direction-aware A$^\ast$ (opt-in) \\
\code{pibt.py} & \ref{sec:pibt} & the PIBT recursion and its rejection rules \\
\code{deadlock.py} & \ref{sec:deadlock} & wait-for cycles, repeated configurations, recovery \\
\code{validate.py} & \ref{sec:pibt-conflicts} & vertex/swap validation, synchronised execution \\
\code{metrics.py} & \ref{sec:metrics} & throughput, service-time percentiles, Jain, runtime \\
\code{simulator.py} & \ref{sec:loop} & the sixteen-step lifelong loop \\
\code{experiments.py} & \ref{sec:results} & ablation, factorial, baseline and hypothesis drivers \\
\code{baselines/} & \ref{sec:baselines} & Token Passing, RHCR, time-expanded A$^\ast$ \\
\code{stats.py} & \ref{sec:statistics} & bootstrap CI and permutation test \\
\bottomrule
\end{tabular}
\caption{Modules and the sections that describe them.}
\label{tab:modules}
\end{table}

\paragraph{Notation.} Throughout: $i, j$ are robot indices and $t$ the current
timestep, advanced by exactly one per \code{Simulator.step()}. Vertices $v, u,
x$ are $(\text{row}, \text{col})$ pairs, and $\pos{i}{t}$ is robot $i$'s
position at time $t$; $g$ and $w$ denote a goal and a waypoint. Greek letters
are tunable weights, all fields of \code{config.Params}, collected in
\Cref{sec:parameters} and glossed in \Cref{sec:glossary}. $\ind{\cdot}$ is the
indicator function, and $\INF$ is the unreachable-distance sentinel.
